% QAI 2026 Overleaf snippets.
% Upload the corresponding PDFs from paper_artifacts/qai2026/figures/.
% Paste table snippets from paper_artifacts/qai2026/tables/ where indicated.

% Place near the beginning of the Framework section.
\begin{figure*}[t]
\centering
\includegraphics[width=0.95\textwidth]{fig_methodology_overview.pdf}
\caption{Overview of the proposed QNN saliency evaluation workflow. Data are standardized and mapped to quantum rotation ranges, passed through either a fully quantum VQC or HQNN, explained with feature-level and parameter-level saliency methods, and evaluated using faithfulness, robustness, stability, and compactness diagnostics.}
\label{fig:methodology-overview}
\end{figure*}

% Place after the model-family definitions in the Framework section.
\begin{figure*}[t]
\centering
\includegraphics[width=0.95\textwidth]{fig_circuit_architectures.pdf}
\caption{Circuit-level architecture summary for the two evaluated model families. The VQC performs feature encoding, entanglement, trainable rotations, and prediction inside the quantum circuit. The HQNN uses the quantum circuit as an expectation-value feature extractor and sends the resulting representation to a classical neural head.}
\label{fig:circuit-architectures}
\end{figure*}

% Place in Results and Discussion -> Experimental Setup.
% Paste table_experimental_setup.tex here.
% Paste table_hyperparameters.tex here.
% Paste table_explanation_hyperparameters.tex here.

% Place near the start of Results and Discussion.
\begin{figure}[t]
\centering
\includegraphics[width=\columnwidth]{fig_accuracy_summary.pdf}
\caption{Cleveland test accuracy across model variants. Iris accuracy is reported separately as a sanity-check benchmark because it is an easier dataset and should not be visually compared against Cleveland.}
\label{fig:cleveland-accuracy}
\end{figure}

% Place in VQC Results.
\begin{figure*}[t]
\centering
\includegraphics[width=0.95\textwidth]{fig_vqc_cleveland_global_importance.pdf}
\caption{Normalized global feature importance for the Cleveland VQC across Integrated Gradients, Grad$\times$Input, noise sensitivity, and occlusion. The dominant rankings consistently emphasize \texttt{thal} and \texttt{cp}.}
\label{fig:vqc-cleveland-global-importance}
\end{figure*}

% Place in VQC Results after the global-importance figure.
\begin{figure}[t]
\centering
\includegraphics[width=\columnwidth]{fig_vqc_cleveland_mean_rank.pdf}
\caption{Consensus Cleveland VQC feature ranking aggregated across attribution methods. Lower mean rank indicates higher consensus importance.}
\label{fig:vqc-cleveland-mean-rank}
\end{figure}

% Place in Faithfulness and Parameter-Level Analysis.
\begin{figure}[t]
\centering
\includegraphics[width=\columnwidth]{fig_vqc_cleveland_deletion_summary.pdf}
\caption{Deletion-based faithfulness summary for the Cleveland VQC. Masking salient features lowers mean true-class confidence, producing an area under the deletion curve (AUDC) of $0.428$.}
\label{fig:vqc-cleveland-deletion}
\end{figure}

% Place in Hybrid QNN Results.
\begin{figure*}[t]
\centering
\includegraphics[width=0.95\textwidth]{fig_hybrid_cleveland_metrics.pdf}
\caption{Hybrid Cleveland explanation metrics for plain and richer encodings. Richer encoding improves rank stability but increases perturbation sensitivity, showing that predictive gains should be interpreted jointly with explanation robustness.}
\label{fig:hybrid-cleveland-metrics}
\end{figure*}

% Place after the main results figures or near the end of Results and Discussion.
% Paste table_results_summary.tex here.
